% Teil II: der durchgerechnete Weg.

\subsection*{Gegeben und gesucht}
Beweisen: $2+4+6+\dots+2n=n(n+1)$ für alle $n\ge1$.

\subsection*{Ansatz}
Vollständige Induktion; im IS den letzten Summanden abspalten.

\subsection*{Induktionsanfang}
$n=1$: links $2$, rechts $1\cdot2=2$. Stimmt.

\subsection*{Induktionsschritt}
Annahme: $2+4+\dots+2n=n(n+1)$. Zu zeigen für $n+1$:
\begin{align}
  2+4+\dots+2n+2(n+1)
  &=n(n+1)+2(n+1) \\
  &=(n+1)(n+2). \tag*{$\checkmark$}
\end{align}
Das ist $n(n+1)$ mit $n$ durch $n+1$ ersetzt. Aus IA und IS folgt die
Behauptung für alle $n\ge1$.

\begin{ergebnis}
\begin{equation*}
  2+4+6+\dots+2n= n(n+1)\quad\text{für alle } n\ge1 .
\end{equation*}
\end{ergebnis}

\begin{kontrolle}
$n=3$: $2+4+6=12$, $3\cdot4=12$; $n=5$: $2+4+6+8+10=30$, $5\cdot6=30$.
Passt.
\end{kontrolle}

\begin{merksatz}
Induktionsschritt: neuen letzten Summanden abspalten, Voraussetzung
einsetzen und den gemeinsamen Faktor $n+1$ ausklammern.
\end{merksatz}
